\section{Архитектура нейросетевой части: MeshCNN и специализированные модели}

\subsection{Общая архитектура}
Решение базируется на архитектуре MeshCNN \cite{hanocka2019meshcnn}, предназначенной для глубокого обучения на полигональных сетках. MeshCNN оперирует рёбрами, используя свёртки, пулинг и анпулинг, которые сохраняют топологию сетки. Основные компоненты:
\begin{itemize}
    \item \textbf{Свёртка на рёбрах}: каждый рёберный признак вычисляется как свёртка признаков двух соседних граней.
    \item \textbf{Пулинг}: объединяет рёбра, уменьшая количество вершин и рёбер, сохраняя при этом топологию.
    \item \textbf{Анпулинг}: восстанавливает исходное число рёбер после пулинга, используется в задачах сегментации.
\end{itemize}

Для каждой операции обучена отдельная модель (классификатор и сегментатор). Входная сетка приводится к фиксированному числу рёбер (2250) и подаётся в MeshCNN. На выходе:
\begin{itemize}
    \item для классификации – вероятность наличия операции в модели;
    \item для сегментации – метка класса для каждого ребра (0 – фон, 1 – операция).
\end{itemize}

\subsection{Функции потерь}
Для классификации используется кросс-энтропийная потеря:
\[
\mathcal{L}_{\text{cls}} = -\frac{1}{N}\sum_{i=1}^{N}\left[y_i\log p_i + (1-y_i)\log(1-p_i)\right],
\]
где $y_i$ – истинная метка (0 или 1), $p_i$ – предсказанная вероятность.

Для сегментации используется взвешенная кросс-энтропия, компенсирующая дисбаланс классов:
\[
\mathcal{L}_{\text{seg}} = -\frac{1}{M}\sum_{j=1}^{M}\sum_{c\in\{0,1\}} w_c \cdot y_{j,c} \log p_{j,c},
\]
где $w_c$ – веса классов (например, $w_1$ выбирается обратно пропорционально частоте класса 1).

\subsection{Обучение}
Обучение проводится на очищенном датасете (21~464 модели) с разбиением 90\% на обучение и 10\% на тест. Используется оптимизатор Adam, начальная скорость обучения 0.001, количество эпох – 50. Контроль качества осуществляется по точности классификации и доле верно классифицированных рёбер.

